\documentclass[10pt,a4paper]{amsart}
\usepackage[margin=19mm]{geometry}
\usepackage{amsmath,amssymb,mathtools}
\usepackage[scr=rsfs,cal=euler]{mathalfa}
\usepackage{xfrac}
\usepackage[unicode,hidelinks,bookmarks=false]{hyperref}
\newcommand{\ie}{i.e.\ }
\newcommand{\cf}{cf.\ }
\newcommand{\resp}{respectively\ }
\newcommand{\Subsection}{\subsection}
\providecommand{\nobreakdash}{\nobreak\hbox{-}}
\setlength{\emergencystretch}{2em}
\multlinegap0pt
\usepackage{bm}
\usepackage{bbm}
\usepackage{dsfont}
\usepackage{xspace}
\usepackage{cancel}
\usepackage{mathtools}
\usepackage{amsthm}
\makeatletter
\@ifundefined{Theorem}{\newtheorem{Theorem}{Theorem}}{}
\@ifundefined{Definition}{\newtheorem{Definition}[Theorem]{Definition}}{}
\@ifundefined{Lemma}{\newtheorem{Lemma}[Theorem]{Lemma}}{}
\makeatother
\title[Periodic $Y$-Systems and Nahm Sums: The Rank 2 Case]{Periodic $Y$-Systems and Nahm Sums: The Rank 2 Case}
\author{Yuma Mizuno}
\date{Source: arXiv:2301.13239v2 (CC BY 4.0)}
\begin{document}
\maketitle

\noindent\emph{Source and licence.} Periodic $Y$-Systems and Nahm Sums: The Rank 2 Case. Version arXiv:2301.13239v2 (CC BY 4.0), distributed under the Creative Commons Attribution 4.0 International licence, as recorded in the arXiv licence metadata for this exact version and shown on the abstract page https://arxiv.org/abs/2301.13239v2. Full rights notice: licenses/arxiv-2301.13239-CC-BY-4.0.txt.\par\smallskip

\noindent\textit{Editorial scope.} Counted unit: the Lemma whose source label is \texttt{lemma:ban list} in the article (source0.tex lines 1140--1188, source label \texttt{lemma:ban list}), delivered together with the article's own complete original proof of it (source0.tex lines 1189--1272, one \texttt{proof} environment of 84 lines). No mathematics is written, repaired or re-derived: statement and proof are the article's own text line for line. Required context retained uncounted: eq:Y1 (lines 159--161); eq:symplectic property (lines 188--195); eq:Y2 (lines 1135--1137). Every definition, display and label of those lines is retained unchanged. Omitted from this package and not redistributed: 1--158, 162--187, 196--1134, 1138--1139, 1273--1969. Nothing inside a retained excerpt is omitted or altered: every retained body is delivered exactly as the article's lines are written, so no omission receipt of any kind accompanies this package. Citations: the retained excerpts carry no citation command at all, so no bibliography entry of the article is transcribed and no reference is redistributed. Reference adaptation: none was needed and none was made; every reference inside the delivered unit resolves inside it, so no unresolved reference or citation is produced. Printed slips kept literal: no printed word, symbol, hypothesis or number of the counted statement or of its proof was altered. Layout and class: the article's own class is \texttt{sigma} with packages including ifpdf, xcolor, eucal, mathtools, bm, cellspace, enumitem, tikz; the wrapper is the lane's pinned \texttt{amsart} editorial wrapper at 10pt on A4, which restates the theorem-like environments the retained text needs and adds the article's own macro definitions that the retained text needs (none), with extra wrapper packages (bm, bbm, dsfont, xspace, cancel, mathtools). The delivered numbering is the unit's own. Assets: no retained span contains a figure environment, a graphics-inclusion command, a PDF-page inclusion, a table or a TikZ picture; no image file is copied into this package, the package declares no asset dependency and no image file is redistributed with it. Licence: version arXiv:2301.13239v2 of this article is distributed under the Creative Commons Attribution 4.0 International licence, as recorded in the arXiv licence metadata for this exact version and shown on the abstract page https://arxiv.org/abs/2301.13239v2. A full rights notice is delivered at licenses/arxiv-2301.13239-CC-BY-4.0.txt. Characters: every retained excerpt is pure ASCII, so the delivered engine, XeLaTeX with its default Latin Modern text font, reports no missing character.
\begin{equation}\label{eq:Y1}
n_{ij;p} = 0\qquad \text{unless} \quad 0 < p < r_i
\end{equation}

\begin{Definition}
 We say that $A_\pm (z) \in \mathbb{Y}_0$ satisfies the \emph{symplectic property} if
 \begin{align}\label{eq:symplectic property}
 A_+ (z) A_-\bigl(z^{-1}\bigr)^{\mathsf{T}} = A_- (z) A_+ \bigl(z^{-1}\bigr)^{\mathsf{T}},
 \end{align}
 where $\mathsf{T}$ is the transpose of a matrix.
 We denote by $\mathbb{Y}$ the subset of $\mathbb{Y}_0$ consisting of pairs satisfying the symplectic property.
\end{Definition}

\begin{equation}\label{eq:Y2}
 n_{ij;p}^{+} = 0 \qquad \text{or} \qquad n_{ij;p}^{-}=0
\end{equation}

\begin{Lemma}\label{lemma:ban list}
 It is impossible that $A_\pm(z) \in \mathbb{Y}$ has the following forms:
 \begin{enumerate}\itemsep=0pt
 \item[$(1)$]
 $A_+(1) =
 \bigl(\begin{smallmatrix}
 2 & -a \\
 * & *
 \end{smallmatrix}\bigr)$,
 $A_-(1) =
 \bigl(\begin{smallmatrix}
 2 & -b \\
 * & *
 \end{smallmatrix}\bigr)$ for odd $a$, $b$.
 \item[$(2)$]
 $A_+(1) =
 \bigl(\begin{smallmatrix}
 2 & -a \\
 * & *
 \end{smallmatrix}\bigr)$,
 $A_-(1) =
 \bigl(\begin{smallmatrix}
 1 & -b \\
 * & *
 \end{smallmatrix}\bigr)$ for odd $a$, $b$.
 \item[$(3)$]
 $A_+(1) =
 \bigl(\begin{smallmatrix}
 1 & -1 \\
 * & *
 \end{smallmatrix}\bigr)$,
 $A_-(1) =
 \bigl(\begin{smallmatrix}
 1 & -1 \\
 * & *
 \end{smallmatrix}\bigr)$.
 \item[$(4)$]
 $A_+(1) =
 \bigl(\begin{smallmatrix}
 1 & 0 \\
 * & *
 \end{smallmatrix}\bigr)$,
 $A_-(1) =
 \bigl(\begin{smallmatrix}
 1 & * \\
 * & *
 \end{smallmatrix}\bigr)$.
 \end{enumerate}
\end{Lemma}

\begin{proof}
 For (2), we can set
 \begin{align*}
 A_+(z) =
 \begin{pmatrix}
 1 + z^r & -f(z) \\
 * & *
 \end{pmatrix},\qquad
 A_-(z) =
 \begin{pmatrix}
 1 + z^r - z^a & -g(z) \\
 * & *
 \end{pmatrix}
 \end{align*}
 for some $f, g \in \mathbb{N}[z]$.
 By the symplectic property \eqref{eq:symplectic property}, we have
 \[
 z^a + z^{a - r} + f(z) g\bigl(z^{-1}\bigr) =
 z^{-a} + z^{r-a} + g(z) f\bigl(z^{-1}\bigr).
 \]
 Since $0 < a$ and $a-r<0$ by \eqref{eq:Y1},
 the sum of the coefficients of the terms in $f(z)g\bigl(z^{-1}\bigr)$ with positive exponents
 is equal to that with negative exponents.
 Since $f(1) g(1)$ (=$ab$) is odd, $f(z) g\bigl(z^{-1}\bigr)$ should contain the constant term $z^0$,
 which contradicts \eqref{eq:Y2}.
 The proof for (1) is similar.

 For (3), we can set
 \begin{align*}
 A_+(z) =
 \begin{pmatrix}
 1 + z^r - z^a & -z^b \\
 * & *
 \end{pmatrix},\qquad
 A_-(z) =
 \begin{pmatrix}
 1 + z^r - z^c & -z^d \\
 * & *
 \end{pmatrix}
 \end{align*}
 with $0< a,b,c,d < r$.
 Without loss of generality, we can assume $a<c$.
 By \eqref{eq:symplectic property}, we have
 \begin{equation*}
 z^{-c} + z^{r-c} + z^a + z^{a-r} + z^{c-a} + z^{d-b}
 = z^{c} + z^{c-r} + z^{-a} + z^{r-a} + z^{a-c} + z^{b-d}.
 \end{equation*}
 Since $c - a > 0$, we see that $c - a$ is equal to $c$, $r-a$, or $b-d$.
 However, the first two cases are impossible by \eqref{eq:Y1}.
 Thus $c - a = b - d$, which implies that
 \begin{equation*}
 z^{-c} + z^{r-c} + z^a + z^{a-r}
 = z^{c} + z^{c-r} + z^{-a} + z^{r-a}.
 \end{equation*}
 Since $a > 0$, we see that $a$ is equal to $c$ or $r-a$.
 However, $a = c$ is impossible by \eqref{eq:Y2}.
 Thus $a = r-a$, which implies that
$
 z^{-c} + z^{r-c}
 = z^{c} + z^{c-r}$.
 Since $c > 0$, we see that $c = r -c$.
 However, this implies that $a = r/2 = c$, which is impossible by \eqref{eq:Y2}.

 For (4), we can set
 \begin{align*}
 A_+(z) =
 \begin{pmatrix}
 1 + z^r -z^a & 0 \\
 * & *
 \end{pmatrix},\qquad
 A_-(z) =
 \begin{pmatrix}
 1 + z^r - z^b & * \\
 * & *
 \end{pmatrix}.
 \end{align*}
 By \eqref{eq:symplectic property}, we have
 \begin{align*}
 z^{a} + z^{a-r} + z^{-b} + z^{r-b} + z^{b-a} =
 z^{-a} + z^{r-a} + z^{b} + z^{b-r} + z^{a-b}.
 \end{align*}
 Comparing the number of the terms with positive and negative exponents,
 we should have $a = b$. This is impossible by \eqref{eq:Y2}.
\end{proof}

\end{document}
